% GENERATED FILE. Edit canonical lessons, then run npm run generate:books.
% canonical-source-hash: fnv1a64:e92c95e281c78fe1
% canonical-lessons: FR-C445-quincaillerie, FR-C445-droguerie, FR-C445-papeterie, FR-C445-mercerie, FR-C445-caviste

\chapter{Petites boutiques}
\label{ch:fr-petites-boutiques}
\hlchaptermodality{445}

\begin{chapteropening}
\textbf{By the end of this chapter:} \emph{I can find small specialist shops: the hardware shop, the household shop, the stationer, the haberdasher and the wine merchant.}

Complete the last lesson of chapter 445: I can find small specialist shops: the hardware shop, the household shop, the stationer, the haberdasher and the wine merchant.
\end{chapteropening}

\section[la quincaillerie]{la quincaillerie — the hardware shop, the ironmonger’s}

\label{lesson:FR-C445-quincaillerie}



\begin{quote}
Before the new one: say the French for the grocer, then the French for the greengrocer.
\end{quote}

\subsection*{You'll want to know: la quincaillerie}
\textbf{la quincaillerie} — \textquotedblleft{}the hardware shop, the ironmonger’s\textquotedblright{}.

\begin{quote}
Il achète des clous à la quincaillerie.
\end{quote}

\begin{quote}
\emph{He buys nails at the hardware shop.}
\end{quote}

\emph{Le quincaillier} runs it.

\begin{cousinweb}
\emph{quincaillerie} is built on \emph{quincaille} 'small metal goods', which is usually traced to an old word that imitates the clink of metal. EtymDB lists it without a source.
\end{cousinweb}

\begin{grammarlens}[title={What you've built}]
One more word for telling, explaining and describing at B1.
\end{grammarlens}

\subsection*{Your turn}
\begin{itemize}
\raggedright
  \item \emph{Say it:} \emph{la quincaillerie}
  \item \emph{Say it:} \emph{la quincaillerie}, once more
  \item \emph{Say it:} a sentence of your own with \emph{la quincaillerie}
  \item \emph{Recall:} say the French for the grocer, then the French for the greengrocer, then say \emph{la quincaillerie} again
\end{itemize}

\begin{tcolorbox}[breakable,colback=teal!4,colframe=teal!35!black,title={Before you move on}]
What does la quincaillerie mean? (\textquotedblleft{}The hardware shop, the ironmonger’s\textquotedblright{}.) Say it once more.
\end{tcolorbox}
\section[la droguerie]{la droguerie — the household goods shop}

\label{lesson:FR-C445-droguerie}



\begin{quote}
Before the new one: say the French for the greengrocer, then the French for the hardware shop.
\end{quote}

\subsection*{You'll want to know: la droguerie}
\textbf{la droguerie} — \textquotedblleft{}the household goods shop\textquotedblright{}.

\begin{quote}
À la droguerie, elle trouve des produits pour le ménage.
\end{quote}

\begin{quote}
\emph{At the household shop she finds cleaning products.}
\end{quote}

Not a chemist's: it sells paint, soap and cleaning products.

\begin{cousinweb}
\emph{droguerie} is built on \emph{drogue}, first 'dried goods, spices'. The origin of \emph{drogue} is not settled; it is sometimes linked to Dutch \emph{droog} 'dry'.
\end{cousinweb}

\begin{grammarlens}[title={What you've built}]
One more word for telling, explaining and describing at B1.
\end{grammarlens}

\subsection*{Your turn}
\begin{itemize}
\raggedright
  \item \emph{Say it:} \emph{la droguerie}
  \item \emph{Say it:} \emph{la droguerie}, once more
  \item \emph{Say it:} a sentence of your own with \emph{la droguerie}
  \item \emph{Recall:} say the French for the greengrocer, then the French for the hardware shop, then say \emph{la droguerie} again
\end{itemize}

\begin{tcolorbox}[breakable,colback=teal!4,colframe=teal!35!black,title={Before you move on}]
What does la droguerie mean? (\textquotedblleft{}The household goods shop\textquotedblright{}.) Say it once more.
\end{tcolorbox}
\section[la papeterie]{la papeterie — the stationer’s}

\label{lesson:FR-C445-papeterie}



\begin{quote}
Before the new one: say the French for the hardware shop, then the French for the household goods shop.
\end{quote}

\subsection*{You'll want to know: la papeterie}
\textbf{la papeterie} — \textquotedblleft{}the stationer’s\textquotedblright{}.

\begin{quote}
Avant la rentrée, la papeterie est pleine d'enfants.
\end{quote}

\begin{quote}
\emph{Before school starts again, the stationer's is full of children.}
\end{quote}

Pronounced \emph{pa-pè-tri}.

\begin{cousinweb}
\emph{papeterie} is built on \emph{papier} 'paper' (EtymDB), which goes back to Latin \emph{papyrus}, from Greek \emph{pápyros}.
\end{cousinweb}

\begin{grammarlens}[title={What you've built}]
One more word for telling, explaining and describing at B1.
\end{grammarlens}

\subsection*{Your turn}
\begin{itemize}
\raggedright
  \item \emph{Say it:} \emph{la papeterie}
  \item \emph{Say it:} \emph{la papeterie}, once more
  \item \emph{Say it:} a sentence of your own with \emph{la papeterie}
  \item \emph{Recall:} say the French for the hardware shop, then the French for the household goods shop, then say \emph{la papeterie} again
\end{itemize}

\begin{tcolorbox}[breakable,colback=teal!4,colframe=teal!35!black,title={Before you move on}]
What does la papeterie mean? (\textquotedblleft{}The stationer’s\textquotedblright{}.) Say it once more.
\end{tcolorbox}
\section[la mercerie]{la mercerie — the haberdasher’s, the sewing shop}

\label{lesson:FR-C445-mercerie}



\begin{quote}
Before the new one: say the French for the household goods shop, then the French for the stationer’s.
\end{quote}

\subsection*{You'll want to know: la mercerie}
\textbf{la mercerie} — \textquotedblleft{}the haberdasher’s, the sewing shop\textquotedblright{}.

\begin{quote}
Ma grand-mère achète ses boutons à la mercerie.
\end{quote}

\begin{quote}
\emph{My grandmother buys her buttons at the haberdasher's.}
\end{quote}

Buttons, thread, ribbons.

\begin{cousinweb}
\emph{mercerie} is built on \emph{mercier} 'dealer in small goods' (EtymDB lists \emph{mercerie}), which goes back to Latin \emph{merx, mercis} 'goods'.
\end{cousinweb}

\begin{grammarlens}[title={What you've built}]
One more word for telling, explaining and describing at B1.
\end{grammarlens}

\subsection*{Your turn}
\begin{itemize}
\raggedright
  \item \emph{Say it:} \emph{la mercerie}
  \item \emph{Say it:} \emph{la mercerie}, once more
  \item \emph{Say it:} a sentence of your own with \emph{la mercerie}
  \item \emph{Recall:} say the French for the household goods shop, then the French for the stationer’s, then say \emph{la mercerie} again
\end{itemize}

\begin{tcolorbox}[breakable,colback=teal!4,colframe=teal!35!black,title={Before you move on}]
What does la mercerie mean? (\textquotedblleft{}The haberdasher’s, the sewing shop\textquotedblright{}.) Say it once more.
\end{tcolorbox}
\section[le caviste]{le caviste — the wine merchant}

\label{lesson:FR-C445-caviste}



\begin{quote}
Before the new one: say the French for the stationer’s, then the French for the haberdasher’s.
\end{quote}

\subsection*{You'll want to know: le caviste}
\textbf{le caviste} — \textquotedblleft{}the wine merchant\textquotedblright{}.

\begin{quote}
Le caviste nous a conseillé un vin de la région.
\end{quote}

\begin{quote}
\emph{The wine merchant recommended a local wine.}
\end{quote}

Built on \emph{la cave} 'cellar'.

\begin{cousinweb}
\emph{caviste} is built on \emph{cave} 'cellar' with \emph{-iste} (EtymDB); \emph{cave} goes back to Latin \emph{cavus} 'hollow'.
\end{cousinweb}

\begin{grammarlens}[title={What you've built}]
One more word for telling, explaining and describing at B1.
\end{grammarlens}

\subsection*{Your turn}
\begin{itemize}
\raggedright
  \item \emph{Say it:} \emph{le caviste}
  \item \emph{Say it:} \emph{le caviste}, once more
  \item \emph{Say it:} a sentence of your own with \emph{le caviste}
  \item \emph{Recall:} say the French for the stationer’s, then the French for the haberdasher’s, then say \emph{le caviste} again
\end{itemize}

\begin{tcolorbox}[breakable,colback=teal!4,colframe=teal!35!black,title={Before you move on}]
What does le caviste mean? (\textquotedblleft{}The wine merchant\textquotedblright{}.) Say it once more.
\end{tcolorbox}
